\subsection{微分算子}

14.1.1. 设 \(k\) 是满足 \(0 \leqslant k \leqslant r\) 的整数，并设 \(\mathbf{P}^k(\mathbf{E})\) 是 E 的 \(k\) 阶截面射流丛（12.6.1）。令

\[
\mathbf {D} ^ {k} (\mathbf {E}, \mathbf {F}) = \mathscr {L} (\mathbf {P} ^ {k} (\mathbf {E}); \mathbf {F}).
\]

它是一个以 X 为底、\(C^{r-k}\) 类的向量丛。这个丛在 X 的开集 U 上的截面等同于一个映射

\[
\mathbf {D}: \mathrm{P} ^ {k} (\mathrm{E}) | \mathrm{U} \rightarrow \mathrm{F} | \mathrm{U}
\]

它与到 U 上的投影交换，且在每个纤维上是线性的；这样的截面称为 U 上的 E \(\to\) F 型、阶 \(\leqslant k\) 的微分算子。设 \(h \in N_{K} \cup \{0\}\)，其中 \(0 \leqslant h \leqslant r-k\)。用 \(\mathcal{D}_{\mathrm{U}}^{k,h}(\mathrm{E}, \mathrm{F})\) 表示 U 上的 E \(\to\) F 型、阶 \(\leqslant k\) 且为 \(C^{h}\) 类（作为 \(\mathrm{D}^{k}(\mathrm{E}, \mathrm{F})\) 的截面，或作为 \(\mathrm{P}^{k}(\mathrm{E})|\mathrm{U}\) 到 \(\mathrm{F}|\mathrm{U}\) 的态射，二者是一回事）的微分算子全体组成的集合。若 \(\mathrm{M} = \mathrm{D}^{k}(\mathrm{E}, \mathrm{F})\)，则有 \(\mathcal{D}_{\mathrm{U}}^{k,h}(\mathrm{E}, \mathrm{F}) = \mathcal{S}_{\mathrm{M}}^{h}(\mathrm{U})\)；它是 \(\mathcal{C}^{h}(\mathrm{U})\) 上的模（参见 7.4.1）。

若 \(0 \leqslant k' \leqslant k\)，则态射 \(r^{k, k'}: \mathrm{P}^{k}(\mathrm{E}) \to \mathrm{P}^{k'}(\mathrm{E})\)（参见 12.6.6）定义了从 \(\mathrm{D}^{k'}(\mathrm{E}, \mathrm{F})\) 到 \(\mathrm{D}^{k}(\mathrm{E}, \mathrm{F})\) 的一个单射；从而每个阶 \(\leqslant k'\) 的微分算子都等同于一个阶 \(\leqslant k\) 的微分算子。一个阶 \(\leqslant k\) 但非阶 \(\leqslant k - 1\) 的微分算子有时称为 \(k\) 阶微分算子。若 \(h \in \mathbf{N}_{\mathbb{K}} \cup \{0\}\) 且 \(h \leqslant r - k\)，且 U 是 X 的开子集，则有：

\[
0 \subset \mathcal {D} _ {\mathrm{U}} ^ {0, h} (\mathrm{E}, \mathrm{F}) \subset \mathcal {D} _ {\mathrm{U}} ^ {1, h} (\mathrm{E}, \mathrm{F}) \subset \dots \subset \mathcal {D} _ {\mathrm{U}} ^ {k, h} (\mathrm{E}, \mathrm{F}).
\]

若 \(r=\infty\) 或 \(\omega\)，用 \(\mathcal{D}_{\mathrm{U}}^{\infty,h}(\mathrm{E},\mathrm{F})\) 表示诸 \(\mathcal{D}_{\mathrm{U}}^{k,h}(\mathrm{E},\mathrm{F})\)（\(k\geqslant0\)）的并；这个并的元素称为 U 上的有界阶微分算子（E \(\rightarrow\) F 型、C \(^{h}\) 类），有时也简称 U 上的微分算子。

14.1.2（“零阶算子”）。我们有 \(\mathbf{P}^0 (\mathbf{E}) = \mathbf{E}\)（12.5.1）以及 \(\mathrm{D}^0 (\mathrm{E},\mathrm{F}) = \mathcal{L}(\mathrm{E};\mathrm{F})\)。若 U 是 X 中的开集，且 \(h\in \mathbf{N}_{\mathbb{K}}\cup \{0\}\)、\(h\leqslant r\)，则 \(\mathcal{D}_{\mathbf{U}}^{0,h}(\mathbf{E},\mathbf{F})\) 的元素是 \(\mathbf{C}^h\) 类的、从 E\textbar U 到 F\textbar U 的态射。

14.1.3（“结构弱化”）。假设 K = R，并设 \(r' \in N_{R}\) 满足 \(r' \leqslant r\)。设 X'、E' 与 F' 分别是由 X、E 与 F 通过弱化结构得到的 \(C^{r'}\) 类流形和向量丛。设 k 是满足 \(0 \leqslant k \leqslant r'\) 的整数。那么 \(P^{k}(E')\) 是一个 \(C^{r'-k}\) 类丛，由 \(P^{k}(E)\) 通过弱化结构得到，对 \(D^{k}(E', F')\) 也有类似的结果。特别地，若 U 是 X 中的开集，且 \(0 \leqslant h \leqslant r - k'\)，则有 \(\mathcal{D}_{\mathrm{U}}^{k,h}(E', F') = \mathcal{D}_{\mathrm{U}}^{k,h}(E, F)\)。

14.1.4. 设 D 是 X 上的一个 E \(\to\) F 型、阶 \(\leqslant k\)（其中 \(0 \leqslant k \leqslant r\)）的微分算子。设 U 是 X 的开子集，s 是 E 在 U 上的一个 \(\mathbf{C}^{m}\) 类截面，其中 \(m \in N _{K} \cup \{0\}\) 且 \(k \leqslant m \leqslant r\)。那么 \(j^{k}(s)\) 是一个 \(\mathbf{C}^{m-k}\) 类截面，即 \(P^{k}(E)\) 在 U 上的截面（12.6.1）；它在 D 下的像记作 \(D_{U}(s)\) 或 D(s)。它是 F 在 U 上的一个截面。假设 D 是 \(\mathbf{C}^{h}\) 类的，其中 \(h=m-k\)。那么 \(D_{U}(s)\) 是 \(\mathbf{C}^{h}\) 类的，且诸映射

\[
\mathrm{D} _ {\mathrm{U}} \colon \mathcal {S} _ {\mathrm{E}} ^ {m} (\mathrm{U}) \rightarrow \mathcal {S} _ {\mathrm{F}} ^ {h} (\mathrm{U})
\]

具有下列性质：

\begin{enumerate}
\def\labelenumi{(\arabic{enumi})}
\item
  \(D_{U}\) 是 K-线性的。
\item
  对每个 \(s \in \mathcal{S}_{\mathbb{E}}^{m}(\mathrm{U})\) 及每个 \(\mathrm{V}\)（\(\mathrm{U}\) 的开集），有 \(\mathrm{D}_{\mathrm{V}}(s|\mathrm{V}) = \mathrm{D}_{\mathrm{U}}(s)|\mathrm{V}\)。
\item
  对每个 \(s \in \mathcal{S}_{\mathbb{E}}^{m}(\mathrm{U})\) 及每个 \(x \in \mathrm{U}\)（满足 \(j_{x}^{k}(s) = 0\)），有 \(D_{\mathrm{U}}(s)(x) = 0\)。
\end{enumerate}

(3') 对 U 内每个坐标系 \(\xi = (\xi^1, \ldots, \xi^n)\) 以及 E 在 U 上的任意标架 \(s = (s_1, \ldots, s_d)\)（7.4.4），存在 F 在 U 上的截面 \(n_{i,\alpha}(1 \leqslant i \leqslant d, \alpha \in \mathbf{N}^n, |\alpha| \leqslant k)\)（\(C^h\) 类的截面），使得对每个族 \((f_i)_{1 \leqslant i \leqslant d}\)（\(\mathcal{C}^m(U)\) 中元素的族），有

\[
\mathrm{D} _ {\mathrm{U}} \left(\sum_ {1 \leqslant i \leqslant d} f _ {i}. s _ {i}\right) = \sum_ {1 \leqslant i \leqslant d, | \alpha | \leqslant k} \Delta_ {\xi} ^ {\alpha} (f _ {i}). n _ {i, \alpha},
\]

\begin{enumerate}
\def\labelenumi{(\arabic{enumi})}
\setcounter{enumi}{3}
\tightlist
\item
  对任意函数 \(f \in \mathcal{C}^m(\mathrm{U})\) 及任意映射 \(\theta\)（从 \(\mathcal{S}_{\mathrm{U}}^{m}(\mathrm{E})\) 到 \(\mathcal{S}_{\mathrm{U}}^{h}(\mathrm{F})\)），用 \(\operatorname{ad}(f)\theta\) 表示映射
\end{enumerate}

\[
s \mapsto f. \theta (s) - \theta (f. s)
\]

（从 \(\mathcal{S}_{\mathrm{U}}^{m}(\mathrm{E})\) 到 \(\mathcal{S}_{\mathrm{U}}^{h}(\mathrm{F})\)）。那么，对每个函数 \(f \in \mathcal{C}^{m}(\mathrm{U})\)，存在一个微分算子 \(\mathbf{L}\)（\(\mathbf{U}\) 上的算子），它是 \(\mathrm{E} \to \mathrm{F}\) 型、阶 \(\leqslant k - 1\) 且 \(\mathbf{C}^{h}\) 类的，使得

\[
(\operatorname{ad} (f) \mathrm{D} _ {\mathrm{U}}) (s) = \mathrm{L} _{\mathrm{U}}(s) \quad \text { for all } s \in \mathscr{S}_{\mathrm{E}}^{m}(\mathrm{U}).
\]

\begin{enumerate}
\def\labelenumi{(\arabic{enumi})}
\setcounter{enumi}{4}
\tightlist
\item
  我们有
\end{enumerate}

\[
\operatorname{ad} (f _ {0}) \dots \operatorname{ad} (f _ {k}) \mathrm{D} _ {\mathrm{U}} = 0 \quad \text { for all } f _ {0}, \dots , f _ {k} \text { in } \mathscr {C} ^ {m} (\mathrm{U}).
\]

若令 I = \{0, . . ., k\}，并对 I 的每个子集 H 令 \(f_{H} = \prod_{i \in H} f_i\)，则前面的关系式等价于：

\[
\sum_ {\mathrm{H} \in \mathrm{I}} (- 1) ^ {\operatorname{Card} (\mathrm{H})} f _ {\mathrm{H}}. \mathrm{D} _ {\mathrm{U}} (f _ {\mathrm{I-H}}. s) = 0
\]

其中 \(f_{0},\ldots,f_{k}\) 取遍 \(\mathcal{C}^{m}(\mathrm{U})\)，\(s\) 取遍 \(\mathcal{S}_{\mathbf{E}}^{m}(\mathrm{U})\)。

14.1.5（“微分算子的刻画”）。设 \(k, m\) 与 \(h = m - k\) 如 14.1.4 所述。对 X 的每个开集 U，设 \(\mathbf{D}_{\mathbf{U}}\) 是从 \(\mathcal{S}_{\mathbf{E}}^{m}(\mathbf{U})\) 到 \(\mathcal{S}_{\mathbf{F}}^{h}(\mathbf{U})\) 的一个映射。假设对 X 的每个开子集 U，14.1.4 的条件 1、2、3（相应地，1、2、3'）都满足。那么存在 \(\mathcal{D}_{\mathbf{X}}^{k,h}(\mathbf{E},\mathbf{F})\) 中唯一的元素 D，使得诸 \(\mathbf{D}_{\mathbf{U}}\) 是相应的映射。\(^{(1)}\) 在下文中，我们把 D 与诸 \(\mathbf{D}_{\mathbf{U}}\) 组成的族恒同。

当 \(m = \infty\) 或 \(\omega\) 时，有 \(h = m\)，并且可以把条件 (3) 与 (3') 换成条件 (4) 或 (5) 之一。

14.1.6（“标量算子”）。假设 E 与 F 都等于平凡丛 \(K_{X}\)。这时 \(E \rightarrow F\) 型的微分算子称为标量微分算子，或简称微分算子；若它是阶 \(\leqslant k\) 的，则它是丛 \(\mathcal{L}(\mathrm{P}^{k}(\mathrm{X}); \mathrm{K}_{\mathrm{X}}) = \mathrm{P}^{k}(\mathrm{X})^{*}\)（丛 \(\mathrm{P}^{k}(\mathrm{X})\) 的对偶）的截面；利用同构

\[
i ^ {- 1} \colon \mathrm{P} ^ {k} (\mathrm{X}) ^ {*} \rightarrow \mathrm{T} ^ {(k)} (\mathrm{X}) \quad (\text { cf. } 13.2.5),
\]

可知阶 \(\leqslant k\) 的标量微分算子等同于向量丛 \(\mathbf{T}^{(k)}(\mathbf{X})\) 的截面，即阶 \(\leqslant k\) 的点分布场。

特别地，取 X 为 \(K^{n}\) 的一个开集，其中 n 是一个 \(\geqslant0\) 的整数。点分布场

\[
\Delta^ {\alpha} \colon x \mapsto \Delta_ {x} ^ {\alpha} \quad (\text { cf. } 1 3. 2. 6)
\]

都是 \(\mathbf{C}^{\omega}\) 类的标量微分算子。对每个 \(h\in\mathbf{N}_{\mathbf{K}}\)，诸 \(\Delta^{\alpha}\)（\((|\alpha|\leqslant k)\)）构成 \(\mathcal{C}^{h}(\mathbf{X})\)-模 \(\mathcal{D}_{\mathbf{X}}^{k,h}(\mathbf{K}_{\mathbf{X}},\mathbf{K}_{\mathbf{X}})\) 的一个基。

14.1.7（“实流形上的复微分算子”）。假设 K = R，且向量丛 E 与 F 都装备了复结构（8.8.1）；设 k 是满足 \(0 \leqslant k \leqslant r\) 的整数。这时丛 \(\mathrm{P}^{k}(\mathrm{E})\) 装备有一个复结构（12.6.3）；丛 \(\mathcal{L}_{\mathbf{C}}(\mathrm{P}^{k}(\mathrm{E}); \mathrm{F})\)（7.8.5，由 \(\mathrm{P}^{k}(\mathrm{E})\) 到 F 的复线性映射组成的丛）是 \(\mathcal{L}(\mathrm{P}^{k}(\mathrm{E}); \mathrm{F}) = \mathrm{D}^{k}(\mathrm{E}, \mathrm{F})\) 的一个向量子丛；我们把它记作 \(\mathrm{D}_{\mathbf{C}}^{k}(\mathrm{E}, \mathrm{F})\)。\(\mathrm{D}_{\mathbf{C}}^{k}(\mathrm{E}, \mathrm{F})\) 的截面称为 \(E \to F\) 型、阶 \(\leqslant k\) 的复微分算子。若 \(h \in N_{K}\) 且 \(h \leqslant r - k\)，且 U 是 X 的开集，我们类似地用 \(\mathcal{D}_{\mathbf{U}}^{k,h}(\mathrm{E}, \mathrm{F})_{\mathbf{C}}\) 表示 \(C^{h}\) 类截面（\(\mathrm{D}_{\mathbf{C}}^{k}(\mathrm{E}, \mathrm{F})\) 在 U 上的截面）所成的空间；它是一个 C-向量空间。本段的其他定义和结论类似地推广到复微分算子；其表述留给读者。

14.1.8（“复合”）。设 G 是一个以 X 为底的 \(C^{r}\) 类向量丛。设 \(k'\) 与 \(k''\) 是正整数，其和 \(k \leqslant r\)，并设 \(D'\)（相应地，\(D''\)）是 X 上的一个 \(E \to F\) 型（相应地，\(F \to G\) 型）、阶 \(\leqslant k'\)（相应地，\(\leqslant k''\)）的微分算子。假设 \(D': P^{k'}(E) \to F\) 是 \(C^{h'}\) 类的，其中 \(h' \in N_{K} \cup \{0\}\) 且 \(k'' \leqslant h' \leqslant r - k'\)；映射

\[
\mathbf {D} _ {*} ^ {\prime} = \mathbf {P} ^ {k ^ {\prime \prime}} (\mathbf {D} ^ {\prime}) \colon \mathbf {P} ^ {k ^ {\prime \prime}} (\mathbf {P} ^ {k ^ {\prime}} (\mathbf {E})) \rightarrow \mathbf {P} ^ {k ^ {\prime \prime}} (\mathbf {F})
\]

该映射是 \(\mathbf{C}^{h'-k''}\) 类的（12.6.3）。设 \(\beta\) 是从 \(\mathbf{P}^{k}(\mathbf{E})\) 到 \(\mathbf{P}^{k''}(\mathbf{P}^{k'}(\mathbf{E}))\) 的典范同态（参见 12.6.5）。把映射

\[
\mathrm{P} ^ {k} (\mathrm{E}) \xrightarrow {\beta} \mathrm{P} ^ {k ^ {\prime \prime}} (\mathrm{P} ^ {k ^ {\prime}} (\mathrm{E})) \xrightarrow {\mathrm{D} _ {*} ^ {\prime}} \mathrm{P} ^ {k ^ {\prime \prime}} (\mathrm{F}) \xrightarrow {\mathrm{D} ^ {\prime \prime}} \mathrm{G},
\]

复合起来，就得到一个 E → G 型、阶 ≤ k 的微分算子。这个算子称为 D'' 与 D' 的复合；记作 D'' ∘ D'。

\(^{1}\) 此外，只需条件 (3′) 对一个坐标系族 \((\xi_{\lambda})\) 和一个定义域覆盖 X 的标架族 \((s_{\lambda})\) 满足即可。

若 U 是 X 的开子集，且 \(s \in \mathcal{S}_{\mathrm{E}}^{k}(\mathrm{U})\)，则有

\[
(\mathbf {D} ^ {\prime \prime} \circ \mathbf {D} ^ {\prime}) (s) = \mathbf {D} ^ {\prime \prime} (\mathbf {D} ^ {\prime} (s))
\]

（其中两边都是 G 在 U 上的截面）；这印证了所采用的术语和记号。

现在假设 \(\mathbf{D}^{\prime \prime}\) 是 \(\mathbf{C}^{h^{\prime \prime}}\) 类的，其中 \(h^{\prime \prime} \in \mathrm{N}_{\mathbb{K}} \cup \{0\}\) 且 \(h^{\prime \prime} \leqslant r-k^{\prime\prime}\)。那么 \(\mathbf{D}^{\prime \prime} \circ \mathbf{D}^{\prime}\) 是 \(\mathbf{C}^{h}\) 类的，其中 \(h=\inf(h'-k'',h'')\)。

假设 E = F = G，且 \(k' \leqslant h''\)，此时 \(D' \circ D''\) 有定义。那么 \(D' \circ D'' - D'' \circ D'\) 是一个 \(E \to E\) 型、阶 \(\leqslant k - 1\) 的微分算子；我们把它记作 \([D', D'']\)。

14.1.9（“结合性”）。沿用 14.1.8 的记号与假设，设 H 是一个 \(\mathbf{C}^r\) 类、以 X 为底的向量丛，并设 \(\mathbf{D}''\) 是 X 上的一个 \(\mathbf{G} \to \mathbf{H}\) 型、阶 \(\leqslant k^{m}\) 的微分算子，其中 \(k' + k'' + k'' \leqslant r\)。我们假设 \(\mathbf{D}'\) 是 \(\mathbf{C}^{h'}\) 类的，其中 \(h' \in \mathrm{N}_{\kappa} \cup \{0\}\) 且 \(k'' + k'' \leqslant h' \leqslant r - k'\)；并假设 \(\mathbf{D}''\) 是 \(\mathbf{C}^{h''}\) 类的，其中 \(h'' \in \mathrm{N}_{\kappa} \cup \{0\}\) 且 \(k'' \leqslant h'' \leqslant r - k''\)。那么复合

\[
\mathbf {D} ^ {\prime \prime \prime} \circ (\mathbf {D} ^ {\prime \prime} \circ \mathbf {D} ^ {\prime}) \quad \text { and } \quad (\mathbf {D} ^ {\prime \prime \prime} \circ \mathbf {D} ^ {\prime \prime}) \circ \mathbf {D} ^ {\prime}
\]

都有定义并且相等。

14.1.10. 例子

\begin{enumerate}
\def\labelenumi{\alph{enumi})}
\tightlist
\item
  设 U 是 X 的开子集，\(D \in \mathcal{D}_{\mathrm{U}}^{k,h}(\mathrm{E}, \mathrm{F})\)，其中 \(h \in N_{K} \cup \{0\}\) 且 \(h \leqslant r-k\)，并设 \(f \in \mathcal{C}^{h+k}(\mathrm{U})\)（相应地，\(f \in \mathcal{C}^{h}(\mathrm{U})\)）。E（相应地，F）中乘以 f 的运算是 U 上的一个阶 \(\leqslant0\)、E \(\to\) E 型（相应地，F \(\to\) F 型）的微分算子（参见 14.1.2）。复合 \(D\circ f\) 与 \(f\circ D\) 有定义且属于 \(\mathcal{D}_{\mathrm{U}}^{k,h}(\mathrm{E},\mathrm{F})\)；我们令
\end{enumerate}

\[
\operatorname{ad} (f) \mathrm{D} = f \circ \mathrm{D} - \mathrm{D} \circ f \in \mathcal {D} _ {\mathrm{U}} ^ {k - 1, h} (\mathrm{E}, \mathrm{F}),
\]

（参见 14.1.4, (4)）。这样就给 \(\mathcal{D}_{\mathrm{U}}^{k,h}(\mathrm{E},\mathrm{F})\) 装备了一个右 \(\mathcal{C}^{h+k}(\mathrm{U})\)-模（相应地，左 \(\mathcal{C}^h (\mathrm{U})\)-模）结构。左模结构与 14.1.1 中的一致。

\begin{enumerate}
\def\labelenumi{\alph{enumi})}
\setcounter{enumi}{1}
\item
  假设 \(r = \infty\) 或 \(r = \omega\)。\(E \to E\) 型、\(C^r\) 类的微分算子构成一个含幺元素的结合 K-代数。
\item
  取 X 为 \(K^{n}\) 的一个开子集。标量微分算子 \(\Delta^{\alpha}\)（14.1.6）满足公式
\end{enumerate}

\[
\Delta^ {\alpha} \circ \Delta^ {\beta} = ((\alpha , \beta)) \Delta^ {\alpha + \beta}.
\]

它们彼此交换。用 \(D_{i}\) 表示算子 \(\Delta^{\varepsilon_{i}}\)（“第 i 个偏导数”）。若 K 的特征为 0，我们有

\[
\Delta^ {\alpha} = \frac {1}{\alpha !} D _ {1} ^ {\alpha_{1}} \dots D _ {n} ^ {\alpha_{n}}.
\]

若 K 的特征为 \(p \neq 0\)，则有 \(\mathbf{D}_{i}^{p} = 0\)（对一切 \(i\)）。

14.1.11（“多微分算子”）。设 \(E_{1}, \ldots, E_{n}\) 是以 X 为底的 \(C^{r}\) 类向量丛，并设 \(k, k_{1}, \ldots, k_{n}\) 是属于 \([0, r]\) 的整数。令

\[
\mathrm{L} (k _ {1}, \dots , k _ {n}) = \mathscr {L} (\mathrm{P} ^ {k _ {1}} (\mathrm{E} _ {1}) \otimes \dots \otimes \mathrm{P} ^ {k _ {n}} (\mathrm{E} _ {n}); \mathrm{F}).
\]

若 \(k_i \leqslant k\)（对一切 \(i\)），则投影 \(r^k, k_i: \mathrm{P}^k(\mathrm{E}_i) \to \mathrm{P}^{k_i}(\mathrm{E}_i)\) 使我们能把丛 \(\mathrm{L}(k_1, \ldots, k_n)\) 等同于丛 \(\mathrm{L}(k) = \mathrm{L}(k, \ldots, k)\) 的一个子丛。存在最小子丛 \(\mathrm{M}(k)\)（\(\mathrm{L}(k)\) 的子丛），它包含所有子丛 \(\mathrm{L}(k_1, \ldots, k_n)\)（\(k_1 + \cdots + k_n = k\)）。截面 \(\mathrm{H}\)（\(\mathrm{M}(k)\) 的截面）称为 \((\mathrm{E}_1, \ldots, \mathrm{E}_n) \to \mathrm{F}\) 型、阶 \(\leqslant k\) 的 n-微分算子（或多微分算子）。若 \(s_i\)（\(1 \leqslant i \leqslant n\)）是 \(\mathbf{C}^m (m \in \mathbb{N}_{\mathbb{K}}, m \geqslant k)\) 类截面（\(\mathbf{E}_i\) 在 X 的开集 U 上的截面），则把 \(\mathrm{H}_{\mathrm{U}}(s_1, \ldots, s_n) = \mathrm{H}(s_1, \ldots, s_n)\) 定义为下述丛的截面 \(j^k(s_1) \otimes \cdots \otimes j^k(s_n)\) 在 H 下的像：

\[
\mathbf {P} ^ {k} (\mathrm{E} _ {1}) \otimes \dots \otimes \mathbf {P} ^ {k} (\mathrm{E} _ {n}).
\]

它是 F\textbar U 的一个截面。若 H 是 \(\mathbf{C}^h\) 类的，则该截面是 \(\mathbf{C}^p\) 类的，其中 \(p = \inf(m - k, h)\)。

当 \(E_{1}, \ldots, E_{n}\) 与 F 都等于 \(K_{x}\) 时，称 H 为标量 n-微分算子（或标量多微分算子）。

例。—— 设 J 是一个有限集，X 是 K \(^{J}\) 的一个开子集。设 H 是 X 上的一个阶 \(\leqslant k\) 的标量 n-微分算子。存在唯一的标量函数族

\[
c (\alpha (1), \dots , \alpha (n)) _ {(\alpha (1), \dots , \alpha (n)) \in \mathbf {N} ^ {\mathrm{J}}}, \quad \sum_ {i = 1} ^ {n} | \alpha (i) | \leqslant k
\]

（定义在 X 上），使得

\[
\mathrm{H} _ {\mathrm{X}} \left(f _ {1}, \dots , f _ {n}\right) = \sum_ {\alpha (1), \dots , \alpha (n)} c (\alpha (1), \dots , \alpha (n)) \Delta^ {\alpha (1)} \left(f _ {1}\right) \dots \Delta^ {\alpha (n)} \left(f _ {n}\right)
\]

对每个函数族 \((f_1, \ldots, f_n)\)（\(\mathbf{C}^k\) 类函数族，\(X\) 上的）成立。要使 \(H\) 是 \(\mathbf{C}^h\) 类的，必须且只需诸函数 \(c(\alpha(1), \ldots, \alpha(n))\) 是 \(\mathbf{C}^h\) 类的。

